% !TeX root = thesis.tex

\chapter{Introduction}
\label{ch:introduction}

Vision-Language Models (VLMs) have demonstrated remarkable capabilities across
a wide range of visual reasoning tasks, yet systematic evaluation across diverse
sensor modalities remains limited. Most existing benchmarks focus exclusively on
standard RGB images or video, leaving open questions about how well these models
generalize to alternative visual representations such as infrared thermal,
event-based, or depth imagery—modalities that are increasingly central to
robotics, autonomous driving, and embodied intelligence applications.

Constructing a high-quality benchmark for multi-sensor VLM evaluation faces two
fundamental challenges. First, manually annotating question--answer (QA) pairs
for multiple aligned sensor modalities is prohibitively expensive and difficult
to scale. Second, automatic annotation pipelines risk producing hallucinated,
unanswerable, or structurally malformed QA items unless rigorous quality control
is applied throughout the generation process.

This report presents an automated pipeline that addresses both challenges.
Starting from multi-sensor video recordings spanning four aligned visual modalities
(RGB, infrared, event, and Marigold depth), the system generates modality-grounded
QA annotations through a structured three-stage prompt chain, filters and validates them with a two-tier quality control mechanism,
and integrates them via a late fusion framework governed by category-driven
reliability profiles.
The resulting benchmark enables controlled evaluation of VLMs under three
distinct protocols: a source-modality baseline, a same-question
modality-replacement experiment, and a native video adapter protocol.

\section{Contributions}
\label{sec:intro_contributions}

The primary contributions of this report are:

\begin{itemize}
    \item An end-to-end automated annotation pipeline that processes four aligned
          visual sensor modalities and produces 5,465 quality-verified QA pairs
          from raw multi-sensor video streams.

    \item A two-tier quality control framework combining rule-based sanitization
          with LLM-assisted semantic verification, achieving a 74.8\% retention rate
          while suppressing hallucinated, unanswerable, and structurally malformed
          annotations.

    \item A late fusion module with category-driven reliability profiles and
          relation-aware cross-modal support scoring that integrates modality-specific
          candidate pools into a coherent, grounded benchmark.

    \item Three controlled VLM benchmarking protocols evaluated across 4B and 8B
          parameter variants of Qwen3-VL, InternVL3.5, and Molmo2, providing
          empirical insights into how model architecture and evidence representation
          influence modality robustness.
\end{itemize}

\section{Report Structure}
\label{sec:intro_structure}

The remainder of this report is organized as follows.
Chapter~\ref{ch:related_work} reviews related work on multimodal datasets,
automated annotation pipelines, and VLM evaluation benchmarks.
Chapter~\ref{ch:approach} describes the proposed annotation and benchmarking
pipeline in detail.
Chapter~\ref{ch:evaluation} presents the experimental setup, benchmark results,
and discussion.
Chapter~\ref{ch:conclusion} summarizes the contributions and outlines directions
for future work.
